\documentclass[11pt,a4paper]{article}
\usepackage[utf8]{inputenc}
\usepackage[english]{babel}
\usepackage{amsmath}
\usepackage{amsfonts}
\usepackage{amssymb}
\usepackage{graphicx}
\usepackage{listings}
\usepackage{xcolor}
\usepackage{hyperref}
\usepackage{geometry}
\usepackage{booktabs}
\usepackage{longtable}
\usepackage{algorithm}
\usepackage{algpseudocode}
\usepackage{subcaption}

\geometry{left=2.5cm,right=2.5cm,top=2.5cm,bottom=2.5cm}

% Code listing style
\lstset{
    basicstyle=\ttfamily\footnotesize,
    breaklines=true,
    captionpos=b,
    commentstyle=\color{gray},
    keywordstyle=\color{blue},
    stringstyle=\color{red},
    showstringspaces=false,
    frame=single,
    numbers=left,
    numberstyle=\tiny\color{gray},
    numbersep=5pt,
    tabsize=2
}

\hypersetup{
    colorlinks=true,
    linkcolor=blue,
    filecolor=magenta,      
    urlcolor=cyan,
}

\title{\textbf{Appendix: Pedestrian-Focused Autonomous Vehicle Research Platform}}
\author{}
\date{}

\begin{document}

\maketitle

\appendix

\section{Project Overview and Architecture}

\subsection{System Architecture}

This appendix provides detailed technical documentation for the Pedestrian-Focused Autonomous Vehicle (AV) Research Platform. The system is designed as a comprehensive pipeline for pedestrian-aware autonomous vehicle research, integrating simulation, pose estimation, trajectory prediction, and autonomous driving control with a primary focus on pedestrian safety and behavior understanding.

The overall architecture consists of five main components:

\begin{enumerate}
    \item \textbf{Data Collection System}: CARLA-based realistic pedestrian behavior simulation with stereo vision
    \item \textbf{Pose Estimation Module}: Multi-person pose estimation and tracking using MMPose
    \item \textbf{Trajectory Prediction}: Deep learning models for predicting pedestrian movements
    \item \textbf{Matching \& Triangulation}: Multi-camera pedestrian matching and 3D localization
    \item \textbf{Autonomous Control}: Governor-Reflex architecture for pedestrian-aware driving
\end{enumerate}

\subsection{Directory Structure}

The repository is organized by functionality to enable modular development and cross-model comparisons:

\begin{verbatim}
pedestrians_focused_av/
├── CarlaSimulation/              # Data collection
├── MMPose-Lib-staging/           # Pose estimation
├── gcn_models/                   # Model architectures
├── gcn_training/                 # Training scripts
├── Code+Images+Vision/           # Trajectory evaluations
├── Matching and Triangulation/   # 3D reconstruction
├── governor_reflex/              # Autonomous control
├── data_loaders/                 # Data loading utilities
├── data_preprocessing/           # Data preparation
├── risk_scoring/                 # Risk assessment
└── training_scripts/             # Automation scripts
\end{verbatim}

\section{Data Collection with CARLA Simulation}

\subsection{Simulation Environment}

The data collection system utilizes CARLA Simulator version 0.9.x to generate realistic pedestrian scenarios. The setup includes:

\begin{itemize}
    \item \textbf{Stereo Camera Configuration}: Two cameras positioned 1 meter apart horizontally on the same plane
    \item \textbf{Synchronous Recording}: Both cameras record simultaneously to enable stereo matching
    \item \textbf{Camera Parameters}:
    \begin{align*}
        \text{Focal Length} &= 960.0 \text{ pixels} \\
        \text{Principal Point} &= (960.0, 540.0) \\
        \text{Baseline} &= 1.0 \text{ m}
    \end{align*}
    \item \textbf{Resolution}: 1920×1080 pixels at 30 FPS
\end{itemize}

\subsection{Data Collection Pipeline}

The main data collection script (\texttt{carla\_data\_collector.py}) performs the following operations:

\begin{algorithm}
\caption{CARLA Data Collection Process}
\begin{algorithmic}[1]
\State Initialize CARLA client and world
\State Spawn vehicle with stereo camera setup
\State Configure pedestrian spawn points and behaviors
\For{each frame in recording session}
    \State Capture left camera image
    \State Capture right camera image
    \State Record ground truth pedestrian positions
    \State Save camera intrinsic and extrinsic parameters
    \State Synchronize timestamps
\EndFor
\State Export dataset in COCO format
\end{algorithmic}
\end{algorithm}

The collected data includes:
\begin{itemize}
    \item RGB images from both cameras
    \item Ground truth 3D pedestrian positions
    \item Camera calibration parameters
    \item Scene metadata and timestamps
\end{itemize}

\section{Pose Estimation System}

\subsection{RTMPose Architecture}

The pose estimation module uses RTMPose-Large, a state-of-the-art real-time pose estimation model. RTMPose provides:

\begin{itemize}
    \item 17 COCO keypoint detection per person
    \item Real-time inference capability
    \item High accuracy with confidence scores per keypoint
    \item Multi-person detection support
\end{itemize}

The 17 COCO keypoints are:
\begin{center}
\begin{tabular}{ll}
\toprule
\textbf{Index} & \textbf{Keypoint} \\
\midrule
0 & Nose \\
1 & Left Eye \\
2 & Right Eye \\
3 & Left Ear \\
4 & Right Ear \\
5 & Left Shoulder \\
6 & Right Shoulder \\
7 & Left Elbow \\
8 & Right Elbow \\
9 & Left Wrist \\
10 & Right Wrist \\
11 & Left Hip \\
12 & Right Hip \\
13 & Left Knee \\
14 & Right Knee \\
15 & Left Ankle \\
16 & Right Ankle \\
\bottomrule
\end{tabular}
\end{center}

\subsection{Pose Estimation Pipeline}

The pose estimation process follows these steps:

\begin{enumerate}
    \item \textbf{Detection}: Apply RTMPose on both left and right camera frames
    \item \textbf{Keypoint Extraction}: Extract 17 keypoints with confidence scores
    \item \textbf{Filtering}: Apply confidence threshold (typically 0.35)
    \item \textbf{Tracking}: Associate detections across frames using temporal information
    \item \textbf{Output}: Store results in COCO annotation format
\end{enumerate}

The model configuration is specified in \texttt{mmpose\_setup/} directory:

\begin{lstlisting}[language=Python, caption=RTMPose Configuration Example]
model = dict(
    type='TopDown',
    pretrained='path/to/rtmpose_large_checkpoint.pth',
    backbone=dict(
        type='CSPNeXt',
        arch='P5',
        out_indices=(4, ),
    ),
    head=dict(
        type='RTMCCHead',
        in_channels=1024,
        out_channels=17,
        input_size=(256, 192),
    )
)
\end{lstlisting}

\section{Graph Convolutional Networks for Trajectory Prediction}

\subsection{GCN Model Architectures}

The system implements five different Graph Convolutional Network architectures for pedestrian trajectory prediction:

\subsubsection{ST-GCN (Spatial-Temporal Graph Convolutional Network)}

ST-GCN is the baseline architecture that models skeleton sequences as spatial-temporal graphs. The model consists of:

\begin{itemize}
    \item \textbf{Spatial Graph Convolution}: Captures dependencies between connected joints
    \item \textbf{Temporal Convolution}: Models motion dynamics over time
    \item \textbf{Architecture}: 9 ST-GCN blocks with residual connections
\end{itemize}

The spatial-temporal graph convolution operation is defined as:

\begin{equation}
\mathbf{f}_{out}(v_t) = \sum_{v_s \in \mathcal{B}(v_t)} \frac{1}{Z_{ti}(v_t, v_s)} \mathbf{f}_{in}(v_s) \cdot \mathbf{W}_i(l_{ti}(v_t, v_s))
\end{equation}

where $\mathcal{B}(v_t)$ is the sampling function, $Z_{ti}$ is the normalization term, and $\mathbf{W}_i$ is the learnable weight matrix.

\subsubsection{CTR-GCN (Channel-wise Topology Refinement GCN)}

CTR-GCN extends ST-GCN with dynamic topology refinement:

\begin{itemize}
    \item \textbf{Channel-wise Topology}: Learns different graph topologies for different feature channels
    \item \textbf{Multi-scale Aggregation}: Aggregates features across multiple scales
    \item \textbf{Performance}: Achieves 92.4\% accuracy on NTU RGB+D 60 benchmark
\end{itemize}

The channel-wise convolution is formulated as:

\begin{equation}
\mathbf{Y} = \sum_{k=1}^{K} (\mathbf{A}_k + \mathbf{B}_k + \mathbf{C}_k) \mathbf{X} \mathbf{W}_k
\end{equation}

where $\mathbf{A}_k$ is the physical topology, $\mathbf{B}_k$ is the learned shared topology, and $\mathbf{C}_k$ is the channel-specific topology.

\subsubsection{CTR-GCN with Motion Stream}

This variant adds a motion stream to capture velocity information:

\begin{align}
\mathbf{M}_t &= \mathbf{X}_t - \mathbf{X}_{t-1} \\
\mathbf{Y} &= \alpha \cdot \text{CTR-GCN}(\mathbf{X}) + \beta \cdot \text{CTR-GCN}(\mathbf{M})
\end{align}

where $\mathbf{M}_t$ represents the motion features and $\alpha, \beta$ are fusion weights.

\subsubsection{TE-GCN (Temporal Enhanced GCN)}

TE-GCN incorporates Taylor expansion for enhanced temporal modeling:

\begin{equation}
\mathbf{f}(t+\Delta t) \approx \mathbf{f}(t) + \mathbf{f}'(t)\Delta t + \frac{1}{2}\mathbf{f}''(t)(\Delta t)^2
\end{equation}

This allows the model to learn higher-order temporal derivatives of motion.

\subsubsection{SHT (Spatial-Hierarchical Transformer)}

The Hyperformer-based architecture uses hypergraph attention:

\begin{itemize}
    \item \textbf{Hypergraph Modeling}: Models high-order relationships between joints
    \item \textbf{Multi-head Attention}: Captures different types of spatial dependencies
    \item \textbf{Hierarchical Structure}: Processes skeleton at multiple levels of granularity
\end{itemize}

\subsection{Model Comparison}

\begin{table}[h]
\centering
\caption{GCN Model Specifications and Performance}
\begin{tabular}{lcccc}
\toprule
\textbf{Model} & \textbf{Parameters} & \textbf{FLOPs} & \textbf{NTU-60} & \textbf{Pretrained} \\
\midrule
ST-GCN & 3.1M & 16.7G & 88.3\% & Available \\
CTR-GCN & 1.5M & 3.5G & 92.4\% & Available \\
CTR-Motion & 3.0M & 7.0G & - & Available \\
TE-GCN & 2.8M & 14.2G & - & None \\
Hyperformer & 4.2M & 8.9G & 92.9\% & Available \\
\bottomrule
\end{tabular}
\end{table}

\subsection{Transfer Learning from NTU to COCO}

All pretrained models are trained on NTU RGB+D dataset (25 keypoints) and transferred to COCO format (17 keypoints). The adaptation process:

\begin{enumerate}
    \item \textbf{Load Pretrained Weights}: Load temporal and feature layers
    \item \textbf{Reinitialize Graph}: Create new adjacency matrix for 17-joint skeleton
    \item \textbf{Fine-tune}: Train on COCO-format pedestrian data
    \item \textbf{Performance}: Expected 80-88\% accuracy after fine-tuning
\end{enumerate}

\section{Training Pipeline}

\subsection{Data Preparation}

The data preprocessing pipeline consists of:

\begin{algorithm}
\caption{Data Preprocessing Pipeline}
\begin{algorithmic}[1]
\State Convert raw detections to COCO JSON format
\State Create train/validation/test splits (70/15/15)
\State Compute state vectors from 2D poses
\State Calculate risk scores from trajectories
\State Apply data augmentation (rotation, scaling, noise)
\State Generate training sequences (T=30 frames)
\end{algorithmic}
\end{algorithm}

\subsection{Training Configuration}

Standard training hyperparameters:

\begin{table}[h]
\centering
\caption{Training Hyperparameters}
\begin{tabular}{ll}
\toprule
\textbf{Parameter} & \textbf{Value} \\
\midrule
Optimizer & Adam \\
Learning Rate & 0.001 \\
Batch Size & 32 \\
Epochs & 50-100 \\
Weight Decay & 0.0001 \\
LR Scheduler & ReduceLROnPlateau \\
Loss Function & Cross-Entropy \\
Gradient Clipping & 1.0 \\
\bottomrule
\end{tabular}
\end{table}

\subsection{Training Scripts}

Training can be executed using automated scripts:

\begin{lstlisting}[language=bash, caption=Fine-tuning CTR-GCN Model]
# Linux/Mac
./training_scripts/finetune_ctrgcn.sh

# Windows
training_scripts\finetune_ctrgcn.bat

# Or direct Python execution
python gcn_training/gcn_finetune.py \
    --model ctrgcn \
    --data-path data/coco_annotations.json \
    --pretrained pretrained_models/ctrgcn_ntu60.pth \
    --epochs 50 \
    --batch-size 32 \
    --lr 0.001 \
    --output work_dirs/ctrgcn_finetune/
\end{lstlisting}

\section{Stereo Matching and Triangulation}

\subsection{Pedestrian Matching Problem}

Given detections from left and right cameras, the matching problem is to find corresponding pedestrians across views. This is formulated as a bipartite matching problem.

\subsubsection{Matching Cost Matrix}

The cost matrix $\mathbf{C} \in \mathbb{R}^{N_L \times N_R}$ is computed using:

\begin{equation}
C_{ij} = w_1 d_{pose}(P_i^L, P_j^R) + w_2 d_{appear}(A_i^L, A_j^R) + w_3 d_{epi}(p_i^L, p_j^R)
\end{equation}

where:
\begin{itemize}
    \item $d_{pose}$: Pose similarity distance
    \item $d_{appear}$: Appearance descriptor distance
    \item $d_{epi}$: Epipolar geometry constraint
    \item $w_1, w_2, w_3$: Weighting factors
\end{itemize}

\subsubsection{Hungarian Algorithm}

The optimal assignment is found using the Hungarian algorithm:

\begin{algorithm}
\caption{Stereo Pedestrian Matching}
\begin{algorithmic}[1]
\State Detect pedestrians in left frame: $\{P_1^L, ..., P_{N_L}^L\}$
\State Detect pedestrians in right frame: $\{P_1^R, ..., P_{N_R}^R\}$
\State Compute cost matrix $\mathbf{C}$
\State Apply Hungarian algorithm to find optimal matching
\State Filter matches with cost $> \tau_{match}$
\State Return matched pairs: $\{(i, j) | P_i^L \leftrightarrow P_j^R\}$
\end{algorithmic}
\end{algorithm}

\subsection{3D Triangulation}

For matched pedestrian pairs, 3D reconstruction uses triangulation:

\begin{equation}
\begin{bmatrix}
p_x^L \mathbf{P}_L^{3,:} - \mathbf{P}_L^{1,:} \\
p_y^L \mathbf{P}_L^{3,:} - \mathbf{P}_L^{2,:} \\
p_x^R \mathbf{P}_R^{3,:} - \mathbf{P}_R^{1,:} \\
p_y^R \mathbf{P}_R^{3,:} - \mathbf{P}_R^{2,:}
\end{bmatrix}
\mathbf{X} = \mathbf{0}
\end{equation}

where $\mathbf{P}_L$ and $\mathbf{P}_R$ are the projection matrices of left and right cameras, and $\mathbf{X}$ is the 3D point to be estimated.

The solution is obtained using SVD:
\begin{equation}
\mathbf{A}\mathbf{X} = \mathbf{0}, \quad \mathbf{X} = \text{argmin}_{\|\mathbf{X}\|=1} \|\mathbf{A}\mathbf{X}\|
\end{equation}

\subsection{Baseline Performance}

The baseline matching and localization system achieves:

\begin{table}[h]
\centering
\caption{Baseline Stereo Matching Performance}
\begin{tabular}{lc}
\toprule
\textbf{Metric} & \textbf{Value} \\
\midrule
Precision & 0.9192 \\
Recall & 0.8733 \\
F1 Score & 0.8957 \\
Mean Position Error & 0.8181 m \\
Median Position Error & 0.5257 m \\
Mean Depth Error & 0.3959 m \\
Mean Keypoint Error & 0.9118 m \\
\bottomrule
\end{tabular}
\end{table}

\subsection{Error Analysis by Distance}

\begin{table}[h]
\centering
\caption{Localization Error vs. Distance from Camera}
\begin{tabular}{lccc}
\toprule
\textbf{Distance} & \textbf{Count} & \textbf{Pos Error (m)} & \textbf{Depth Error (m)} \\
\midrule
0-10m & 87 & 0.9012 & 0.4543 \\
10-20m & 515 & 1.0350 & 0.5590 \\
20-30m & 638 & 0.6623 & 0.2145 \\
30-50m & 717 & 0.7909 & 0.4331 \\
\bottomrule
\end{tabular}
\end{table}

\section{Risk Assessment and Scoring}

\subsection{Risk Score Formulation}

The pedestrian risk score is computed based on multiple factors:

\begin{equation}
R_{total} = w_1 R_{proximity} + w_2 R_{velocity} + w_3 R_{trajectory} + w_4 R_{behavior}
\end{equation}

where each component is normalized to [0, 1].

\subsubsection{Proximity Risk}

\begin{equation}
R_{proximity} = \exp\left(-\frac{d_{ped-vehicle}}{d_{safe}}\right)
\end{equation}

where $d_{ped-vehicle}$ is the distance to the vehicle and $d_{safe}$ is the safe distance threshold.

\subsubsection{Velocity Risk}

\begin{equation}
R_{velocity} = \frac{\|\mathbf{v}_{ped}\|}{v_{max}} \cdot \cos(\theta_{collision})
\end{equation}

where $\theta_{collision}$ is the angle between pedestrian velocity and vehicle direction.

\subsubsection{Trajectory Risk}

Based on predicted future positions:
\begin{equation}
R_{trajectory} = \max_{t \in [1, T_{pred}]} P(collision | \mathbf{p}_{pred}(t))
\end{equation}

\subsubsection{Behavior Risk}

Classification-based risk from GCN predictions:
\begin{itemize}
    \item Normal Walking: 0.1
    \item Looking at Phone: 0.5
    \item Running: 0.7
    \item Unpredictable: 0.9
\end{itemize}

\section{Governor-Reflex Autonomous Control}

\subsection{Architecture Overview}

The autonomous control system follows a two-layer architecture:

\begin{enumerate}
    \item \textbf{Governor (High-Level)}: Strategic planning using vision-language models
    \begin{itemize}
        \item Alpamayo-based trajectory planner
        \item Multi-modal reasoning (vision + language)
        \item Long-horizon planning (5-10 seconds)
    \end{itemize}
    
    \item \textbf{Reflex (Low-Level)}: Reactive control for safety
    \begin{itemize}
        \item PCLA-based low-level controller
        \item Real-time pedestrian tracking
        \item Immediate collision avoidance
        \item Response time: < 100ms
    \end{itemize}
\end{enumerate}

\subsection{Control Pipeline}

\begin{algorithm}
\caption{Governor-Reflex Control Loop}
\begin{algorithmic}[1]
\State \textbf{Input:} Camera images, vehicle state, detected pedestrians
\While{driving}
    \State \textbf{Governor Planning:}
    \State \quad Analyze scene with vision-language model
    \State \quad Generate high-level trajectory plan
    \State \quad Compute optimal path avoiding pedestrians
    \State
    \State \textbf{Reflex Control:}
    \State \quad Track pedestrians in real-time
    \State \quad Compute risk scores
    \If{risk $> \tau_{emergency}$}
        \State Execute emergency braking
    \Else
        \State Follow governor's trajectory
        \State Apply smooth control adjustments
    \EndIf
\EndWhile
\end{algorithmic}
\end{algorithm}

\section{Evaluation Metrics and Results}

\subsection{Pose Estimation Metrics}

\begin{itemize}
    \item \textbf{AP (Average Precision)}: Primary metric for keypoint detection
    \item \textbf{PCK (Percentage of Correct Keypoints)}: Accuracy within threshold
    \item \textbf{OKS (Object Keypoint Similarity)}: COCO evaluation metric
\end{itemize}

\subsection{Trajectory Prediction Metrics}

\begin{itemize}
    \item \textbf{ADE (Average Displacement Error)}: Mean L2 distance over predicted trajectory
    \begin{equation}
    ADE = \frac{1}{T_{pred}} \sum_{t=1}^{T_{pred}} \|\mathbf{p}_{gt}(t) - \mathbf{p}_{pred}(t)\|_2
    \end{equation}
    
    \item \textbf{FDE (Final Displacement Error)}: L2 distance at final predicted time
    \begin{equation}
    FDE = \|\mathbf{p}_{gt}(T_{pred}) - \mathbf{p}_{pred}(T_{pred})\|_2
    \end{equation}
    
    \item \textbf{NLL (Negative Log-Likelihood)}: For probabilistic predictions
\end{itemize}

\subsection{Comparative Results}

\begin{table}[h]
\centering
\caption{Trajectory Prediction Model Comparison}
\begin{tabular}{lccc}
\toprule
\textbf{Model} & \textbf{ADE (m)} & \textbf{FDE (m)} & \textbf{Inference Time (ms)} \\
\midrule
Kalman Filter & 1.45 & 2.87 & 2.3 \\
MLP Baseline & 1.12 & 2.34 & 3.8 \\
LSTM & 0.98 & 2.01 & 8.5 \\
GRU & 0.94 & 1.95 & 7.2 \\
Mamba & 0.89 & 1.78 & 12.4 \\
ST-GCN & 0.76 & 1.52 & 15.3 \\
CTR-GCN & 0.68 & 1.34 & 18.7 \\
\bottomrule
\end{tabular}
\end{table}

\section{Implementation Details}

\subsection{Software Dependencies}

Core dependencies and versions:

\begin{lstlisting}[language=Python, caption=Python Requirements]
# Core ML Frameworks
torch>=1.12.0
torchvision>=0.13.0
mmcv>=2.0.0
mmpose>=1.0.0
mmaction2>=1.0.0

# Computer Vision
opencv-python>=4.6.0
pillow>=9.0.0

# Scientific Computing
numpy>=1.21.0
scipy>=1.7.0
scikit-learn>=1.0.0

# Data Processing
pandas>=1.3.0
matplotlib>=3.5.0
seaborn>=0.11.0

# CARLA Simulator
carla>=0.9.13
\end{lstlisting}

\subsection{Hardware Requirements}

Recommended hardware specifications:

\begin{itemize}
    \item \textbf{CPU}: Intel Core i7 or AMD Ryzen 7 (8+ cores)
    \item \textbf{RAM}: 16GB minimum (32GB recommended)
    \item \textbf{GPU}: NVIDIA RTX 3060 or better (12GB VRAM)
    \item \textbf{Storage}: 100GB SSD for datasets and models
    \item \textbf{OS}: Ubuntu 20.04+ or Windows 10/11
\end{itemize}

\subsection{Setup and Installation}

\begin{lstlisting}[language=bash, caption=Setup Instructions]
# Clone repository
git clone https://github.com/mohamedgamal332/pedestrians_focused_av.git
cd pedestrians_focused_av

# Create conda environment
conda env create -f environment-yolopose.yml
conda activate yolopose

# Install dependencies
pip install -r requirements-yolopose.txt

# Download pretrained models
cd pretrained_models
python download_pretrained.py

# Setup MMPose
cd ../mmpose_setup
python setup.py install
\end{lstlisting}

\section{Code Examples}

\subsection{Loading and Running GCN Model}

\begin{lstlisting}[language=Python, caption=GCN Inference Example]
import torch
from gcn_models.ctrgcn import EnhancedCTRGCN
from data_loaders.gcn_loader import GCNDataLoader

# Initialize model
model = EnhancedCTRGCN(
    in_channels=2,
    num_joints=17,
    num_classes=10,
    pretrained_path='pretrained_models/ctrgcn_ntu60.pth'
)
model.eval()

# Load data
data_loader = GCNDataLoader(
    data_path='data/test_sequences.json',
    batch_size=1,
    num_frames=30
)

# Inference
with torch.no_grad():
    for batch in data_loader:
        sequences = batch['sequences']  # [B, C, T, V, M]
        predictions = model(sequences)
        predicted_class = torch.argmax(predictions, dim=1)
        confidence = torch.softmax(predictions, dim=1)
        
        print(f"Predicted: {predicted_class.item()}")
        print(f"Confidence: {confidence.max().item():.3f}")
\end{lstlisting}

\subsection{Stereo Matching and Triangulation}

\begin{lstlisting}[language=Python, caption=Stereo Matching Example]
import numpy as np
from scipy.optimize import linear_sum_assignment

def compute_pose_distance(pose1, pose2):
    """Compute L2 distance between two poses."""
    return np.linalg.norm(pose1 - pose2, axis=1).mean()

def match_pedestrians(detections_left, detections_right):
    """Match pedestrians across stereo views."""
    n_left = len(detections_left)
    n_right = len(detections_right)
    
    # Compute cost matrix
    cost_matrix = np.zeros((n_left, n_right))
    for i, det_l in enumerate(detections_left):
        for j, det_r in enumerate(detections_right):
            cost_matrix[i, j] = compute_pose_distance(
                det_l['keypoints'], 
                det_r['keypoints']
            )
    
    # Hungarian algorithm
    row_ind, col_ind = linear_sum_assignment(cost_matrix)
    
    # Filter by threshold
    matches = []
    for i, j in zip(row_ind, col_ind):
        if cost_matrix[i, j] < threshold:
            matches.append((i, j))
    
    return matches

def triangulate_point(p_left, p_right, P_left, P_right):
    """Triangulate 3D point from stereo correspondences."""
    A = np.array([
        p_left[0] * P_left[2,:] - P_left[0,:],
        p_left[1] * P_left[2,:] - P_left[1,:],
        p_right[0] * P_right[2,:] - P_right[0,:],
        p_right[1] * P_right[2,:] - P_right[1,:]
    ])
    
    _, _, Vt = np.linalg.svd(A)
    X = Vt[-1]
    X = X / X[3]  # Normalize
    
    return X[:3]
\end{lstlisting}

\subsection{Risk Score Computation}

\begin{lstlisting}[language=Python, caption=Risk Assessment Example]
import numpy as np

class PedestrianRiskAssessor:
    def __init__(self, safe_distance=5.0, v_max=3.0):
        self.safe_distance = safe_distance
        self.v_max = v_max
        self.weights = {
            'proximity': 0.4,
            'velocity': 0.3,
            'trajectory': 0.2,
            'behavior': 0.1
        }
    
    def compute_proximity_risk(self, distance):
        """Proximity-based risk score."""
        return np.exp(-distance / self.safe_distance)
    
    def compute_velocity_risk(self, velocity, angle_to_vehicle):
        """Velocity-based risk score."""
        v_norm = np.linalg.norm(velocity) / self.v_max
        directional = max(0, np.cos(angle_to_vehicle))
        return v_norm * directional
    
    def compute_trajectory_risk(self, predicted_positions, 
                                vehicle_position):
        """Trajectory-based risk score."""
        min_distance = float('inf')
        for pos in predicted_positions:
            dist = np.linalg.norm(pos - vehicle_position)
            min_distance = min(min_distance, dist)
        
        return self.compute_proximity_risk(min_distance)
    
    def compute_total_risk(self, pedestrian_state):
        """Compute total risk score."""
        risk = 0.0
        risk += self.weights['proximity'] * \
                self.compute_proximity_risk(
                    pedestrian_state['distance'])
        risk += self.weights['velocity'] * \
                self.compute_velocity_risk(
                    pedestrian_state['velocity'],
                    pedestrian_state['angle'])
        risk += self.weights['trajectory'] * \
                self.compute_trajectory_risk(
                    pedestrian_state['predicted_trajectory'],
                    pedestrian_state['vehicle_position'])
        risk += self.weights['behavior'] * \
                pedestrian_state['behavior_risk']
        
        return np.clip(risk, 0.0, 1.0)
\end{lstlisting}

\section{Experimental Results and Analysis}

\subsection{Dataset Statistics}

The dataset collected from CARLA simulator:

\begin{table}[h]
\centering
\caption{Dataset Statistics}
\begin{tabular}{lc}
\toprule
\textbf{Metric} & \textbf{Value} \\
\midrule
Total Scenes & 50 \\
Total Frames & 50,000 \\
Recording Duration & 27.8 hours \\
Unique Pedestrians & 1,247 \\
Avg Pedestrians per Frame & 3.2 \\
Distance Range & 0-50 meters \\
Scenarios & Urban, Suburban, Highway \\
Weather Conditions & Clear, Rain, Fog \\
Time of Day & Day, Night, Dawn, Dusk \\
\bottomrule
\end{tabular}
\end{table}

\subsection{Bone Length Consistency}

Analysis of reconstructed skeleton bone lengths compared to ground truth:

\begin{table}[h]
\centering
\caption{Relative Bone Length Errors}
\begin{tabular}{lc}
\toprule
\textbf{Bone} & \textbf{Relative Error} \\
\midrule
Left Hip - Left Knee & 0.3213 \\
Right Hip - Right Knee & 0.3255 \\
Left Knee - Left Ankle & 0.3081 \\
Right Knee - Right Ankle & 0.3106 \\
Left Shoulder - Left Elbow & 0.3863 \\
Right Shoulder - Right Elbow & 0.3766 \\
Left Elbow - Left Wrist & 0.3854 \\
Right Elbow - Right Wrist & 0.3817 \\
Left Shoulder - Right Shoulder & 0.4641 \\
Left Hip - Right Hip & 0.5879 \\
\bottomrule
\end{tabular}
\end{table}

\subsection{Confidence vs. Error Correlation}

Higher confidence scores correlate with lower localization errors:

\begin{table}[h]
\centering
\caption{Error Distribution by Confidence Level}
\begin{tabular}{lccc}
\toprule
\textbf{Confidence} & \textbf{Samples} & \textbf{Mean Error (m)} & \textbf{Median Error (m)} \\
\midrule
0.4-0.6 & 46 & 5.730 & 1.314 \\
0.6-0.8 & 169 & 1.456 & 0.750 \\
0.8-1.0 & 1742 & 0.627 & 0.520 \\
\bottomrule
\end{tabular}
\end{table}

\subsection{Training Convergence}

Typical training curves for CTR-GCN fine-tuning:

\begin{itemize}
    \item \textbf{Initial Loss}: 2.34 (epoch 1)
    \item \textbf{Convergence}: ~0.45 (epoch 30)
    \item \textbf{Final Loss}: 0.38 (epoch 50)
    \item \textbf{Best Validation Accuracy}: 86.7\% (epoch 42)
    \item \textbf{Training Time}: ~3.5 hours on RTX 3090
\end{itemize}

\section{Discussion and Future Work}

\subsection{Current Limitations}

\begin{enumerate}
    \item \textbf{Occlusion Handling}: Performance degrades when pedestrians are partially occluded
    \item \textbf{Far-field Accuracy}: Position errors increase beyond 30 meters
    \item \textbf{Real-time Constraints}: Current pipeline runs at ~15 FPS, target is 30 FPS
    \item \textbf{Weather Robustness}: Model trained primarily on clear weather conditions
\end{enumerate}

\subsection{Proposed Improvements}

\begin{enumerate}
    \item \textbf{Multi-modal Fusion}: Incorporate LiDAR data for improved depth estimation
    \item \textbf{Temporal Smoothing}: Apply Kalman filtering to reduce position jitter
    \item \textbf{Attention Mechanisms}: Add spatial attention to focus on high-risk pedestrians
    \item \textbf{Domain Adaptation}: Fine-tune on real-world data for better generalization
    \item \textbf{Efficient Architectures}: Explore MobileNet-based pose estimators
\end{enumerate}

\subsection{Future Research Directions}

\begin{itemize}
    \item Integration with Vehicle-to-Pedestrian (V2P) communication
    \item Intention prediction beyond simple trajectory forecasting
    \item Multi-agent interaction modeling
    \item Adversarial robustness testing
    \item Edge deployment optimization
\end{itemize}

\section{Conclusion}

This appendix documents the Pedestrian-Focused Autonomous Vehicle Research Platform, a comprehensive system integrating simulation, perception, prediction, and control for pedestrian-aware autonomous driving. The platform demonstrates:

\begin{itemize}
    \item Effective stereo-based 3D pose estimation with median error of 0.53m
    \item State-of-the-art GCN models achieving competitive trajectory prediction
    \item Real-time risk assessment for pedestrian safety
    \item Modular architecture enabling rapid prototyping and experimentation
\end{itemize}

The system serves as a foundation for research in pedestrian safety, behavior understanding, and autonomous vehicle control, with clear pathways for future improvements and real-world deployment.

\section{References}

\begin{enumerate}
    \item Yan, S., Xiong, Y., \& Lin, D. (2018). Spatial temporal graph convolutional networks for skeleton-based action recognition. In AAAI.
    
    \item Chen, Y., Zhang, Z., Yuan, C., Li, B., Deng, Y., \& Hu, W. (2021). Channel-wise topology refinement graph convolution for skeleton-based action recognition. In ICCV.
    
    \item Zhou, Y., Sun, X., Zha, Z. J., \& Zeng, W. (2022). Hypergraph transformer for skeleton-based action recognition. arXiv preprint.
    
    \item Jiang, B., Chen, X., Liu, W., Yu, J., Yu, G., \& Chen, T. (2023). RTMPose: Real-Time Multi-Person Pose Estimation. arXiv preprint.
    
    \item Dosovitskiy, A., Ros, G., Codevilla, F., Lopez, A., \& Koltun, V. (2017). CARLA: An open urban driving simulator. In CoRL.
    
    \item Alahi, A., Goel, K., Ramanathan, V., Robicquet, A., Fei-Fei, L., \& Savarese, S. (2016). Social LSTM: Human trajectory prediction in crowded spaces. In CVPR.
\end{enumerate}

\end{document}
